% ============================================================================
\chapter{논의}\label{ch:discussion}
% ============================================================================

\section{가설 검증 요약}\label{sec:disc_summary}

\tabref{tab:hypotheses}는 \secref{sec:hypothesis}에서 세운 세 가설의 검증 결과이다.

\begin{table}[htbp]
\centering\small
\caption{가설 검증 요약.}
\label{tab:hypotheses}
\begin{tabular}{p{0.1\linewidth}p{0.47\linewidth}p{0.14\linewidth}p{0.18\linewidth}}
\toprule
가설 & 핵심 근거 & 판정 & 남은 단서 \\
\midrule
H1 & nano Mask \mAP{} 0.6272, small 0.6239. 차이는 fold 표준편차 이내이고 small은 3.5배 크다. & 지지 & -- \\
H2 & 영상 단위 검정에서 Dice·BIoU·HD95 모두 유의($p\le0.014$). 추론 GFLOPs 동일(9.73). & 지지 &
  개선 폭은 계수 선택 편향 가능성(\secref{sec:sweep}), 교차 데이터셋은 판정 불가 \\
H3 & OpenVINO INT8: 3.71 MB, CPU 4.10 ms(15.3배), Dice 손실 0.001. & 지지 (OpenVINO) &
  ONNX Runtime INT8은 기준 미달 \\
\bottomrule
\end{tabular}
\end{table}

\section{mAP가 경계 개선을 반영하지 못하는 이유}\label{sec:disc_map}

제안 모델은 경계 지표를 유의하게 개선했지만 Mask \mAP{}는 0.6272에서 0.6267로 거의 변하지 않았다.
이는 지표의 구조로 설명할 수 있다. \secref{sec:disc_lambda}에서 계산하듯 경계 밴드는 결절 상자 안 픽셀의 약 16\%에
불과하므로, 경계 부근이 몇 픽셀 개선되어도 마스크 IoU는 조금밖에 변하지 않는다. mAP는 이 IoU가 0.50, 0.55, \dots, 0.95의
임계값을 넘는지만 세므로, IoU가 임계값을 넘나들지 않는 한 개선이 점수에 나타나지 않는다. 게다가 mAP는 검출 순위와
상자 품질에도 크게 좌우된다.

반면 HD95는 경계의 최악 오차를, Boundary IoU는 경계 띠에서의 겹침만을 재므로 같은 변화에 훨씬 민감하다.
\secref{sec:exp_sig}에서 개선이 일부 영상의 큰 개선에서 나온다는 점을 보였는데, 바로 그런 영상, 즉 베이스라인의 경계가
음영 구간으로 새어 나가던 영상에서 HD95가 크게 줄었다. 이 결과는 \secref{sec:problem}의 P1, 즉 영역 지표만으로는
경계 품질을 판정할 수 없다는 문제 정의를 실험적으로 뒷받침한다.

\section{경계 강조 계수의 해석}\label{sec:disc_lambda}

스윕에서 $\lambda=2$는 $\lambda=10$보다 HD95에서 일관되게 나았다. 경계를 더 강하게 강조할수록 경계가 좋아지리라는 직관과
반대되는 결과이며, \eqnref{eq:weight}의 가중치가 정규화되지 않는다는 점에서 원인을 찾을 수 있다.

가중치 $w_p = 1 + \lambda\mathcal{B}_p$는 모든 픽셀에서 1 이상이므로, $\lambda$를 키우면 경계 픽셀의 상대적 비중이 커지는 동시에
마스크 손실 전체의 크기도 커진다. 정답 상자 안에서 경계 밴드가 차지하는 비율을 $\rho$라 하면 상자 안 평균 가중치는
$\mathbb{E}[w] = 1 + \lambda\rho$이다. 학습 데이터의 결절 3,767개를 640 레터박스 후 마스크 해상도(1/4)로 래스터화하여 계산한
$\rho$의 중앙값은 TN3K 16.3\%, DDTI 17.0\%였다. 이에 따르면 $\lambda=2$에서 마스크 손실은 평균 약 1.33배, $\lambda=10$에서는
약 2.64배(사분위 범위 2.09--3.46배)로 커진다. 학습 중 모자이크와 스케일 증강으로 결절의 크기가 바뀌므로 정확한 값은 다르지만,
$\lambda=10$이 마스크 손실을 두 배 이상 키운다는 점은 달라지지 않는다.

즉 $\lambda$는 '경계 강조'와 '분할 손실의 전체 가중치'라는 두 가지를 함께 바꾼다. $\lambda=10$에서는 분할 손실이 상자·분류
손실보다 과도하게 커져 \eqnref{eq:total_loss}의 손실 균형이 흐트러지고, 이것이 Mask \mAP{} 하락과 경계 개선 폭의 감소로
나타난 것으로 해석된다. 초기값 $\lambda=5$에서 mAP가 떨어졌던 것(0.627$\rightarrow$0.621)도 같은 맥락이다.
후속 연구에서는 $w$를 $\mathbb{E}[w]$로 나누어 손실 크기를 고정하고 경계의 상대 비중만 바꾸는 정규화를 적용하면,
두 효과를 분리하여 $\lambda$를 더 넓은 범위에서 탐색할 수 있을 것이다.

\section{교차 데이터셋 결과의 해석}\label{sec:disc_lodo}

\subsection{성능이 낮은 이유: 데이터의 차이}

LODO에서는 베이스라인과 제안 모델 모두 in-domain보다 성능이 크게 낮았다(Mask \mAP{} 0.30--0.45 대 0.63). 특히
DDTI$\rightarrow$TN3K 방향이 낮았는데, \secref{sec:ds_gap}의 분석이 그 이유를 설명한다. DDTI는 637장으로 작고,
모든 영상이 단일 결절이며, 결절 크기가 영상의 1.5--24\%에 몰려 있다. 이 데이터만으로 학습한 모델은 평가 대상인 TN3K 전체 3,493장 중
8.4\%(294장)를 차지하는 다중 결절 영상과, 영상의 1\% 미만인 작은 결절 또는 40\%를 넘는 큰 결절을 학습한 적이 없다.
영상 형식도 다르다. DDTI는 장비 화면 전체를 담고 있어 결절 주변에 UI 요소가 함께 있는 반면, TN3K는 결절 주변을
잘라낸 확대 영상이다. 따라서 LODO의 낮은 성능은 모델 구조의 결함이라기보다 학습 데이터가 평가 분포를 포함하지
못한 결과로 보는 것이 타당하다.

\subsection{경계 정교화의 효과가 사라진 이유}

in-domain에서 유의했던 경계 개선이 LODO에서 사라진 데에는 두 가지 해석이 가능하다. 하나는 경계 가중 손실이
학습 데이터의 주석 경계에 더 강하게 맞추도록 하므로, 주석 방식이 다른 데이터셋에서는 그 이득이 줄어든다는 것이다.
TN3K$\rightarrow$DDTI의 Boundary IoU가 세 시드 모두에서 제안 모델이 낮았던 점은 이 해석과 어울린다. 다른 하나는
도메인 이동으로 결절 자체를 놓치거나 엉뚱한 영역을 잡는 큰 오류가 늘어나, 경계 몇 픽셀의 개선이 그 오류에 묻혔다는
것이다. LODO에서 HD95가 20\%를 넘는다는 것은 경계 정밀도보다 검출 자체의 실패가 지배적이라는 뜻이다.
현재의 실험으로는 두 해석을 구분할 수 없다.

\subsection{단일 실행으로 결론을 내리면 안 되는 이유}

본 연구의 교차 데이터셋 결과는 실행 간 변동의 중요성을 잘 보여 준다. 시드 42 하나만 보면 TN3K$\rightarrow$DDTI에서
제안 모델의 Dice가 0.118 떨어졌고, 637장 영상 단위 검정의 $p$값은 $10^{-10}$보다 작았다. 이 결과만으로 논문을 썼다면
"경계 정교화는 교차 도메인에서 성능을 크게 해친다"는 강한 결론을 내렸을 것이다. 그러나 시드 두 개를 더하자 그 차이는
대부분 시드 42의 제안 모델 한 번의 실행에서 나온 것으로 드러났다.

여기서 영상 단위 검정과 시드 단위 검정이 답하는 질문이 다르다는 점이 중요하다. 영상 단위 검정은 이미 학습된 특정 모델들을
고정하고 테스트 영상의 표본 변동만 오차로 삼으므로, "이 두 모델 중 어느 쪽이 새 영상에서 나은가"에 답한다. 학습을 다시 하면
달라질 수 있는 부분, 즉 학습의 무작위성은 오차에 포함되지 않는다. 반면 "이 학습 방법이 다른 방법보다 나은가"라는 기법에
관한 질문은 학습을 반복했을 때의 변동을 오차로 삼아야 하며, 이를 반영하는 것이 시드 단위 검정이다\citep{bouthillier2021variance}.
in-domain 결과에서도 같은 구분이 적용된다. 영상 단위 검정이 유의하고 5개 fold 모두에서 Dice가 개선되었다는 점은 강한 근거이지만,
fold 단위 검정은 검정력 한계로 유의수준에 이르지 못했으므로(\secref{sec:exp_sig}), 기법의 효과를 확정하려면 더 많은 반복이
도움이 된다.

\section{양자화 결과의 해석}\label{sec:disc_quant}

두 INT8 경로의 CPU 지연이 44.00 ms와 4.10 ms로 10배 이상 차이 난 가장 큰 원인은 양자화에서 제외된 연산의 범위이다.
제안 모델의 추론 그래프에서 합성곱 연산량을 층별로 집계하면, 백본이 31.8\%, 넥이 16.4\%, 분할 헤드가 51.8\%를 차지한다.
헤드의 비중이 큰 이유는 프로토타입 경로가 입력의 1/4 해상도(160$\times$160)에서 64채널 3$\times$3 합성곱을 수행하기
때문으로, 프로토타입 경로 하나가 전체의 27.9\%이다.

ONNX Runtime 경로는 출력 붕괴를 막기 위해 헤드 전체를 부동소수점으로 남겼으므로, 연산량의 절반 이상이 양자화되지
않았다. 양자화된 나머지 48\%가 아무리 빨라져도 전체 지연은 헤드의 부동소수점 연산 시간 아래로 내려갈 수 없다.
여기에 양자화된 구간과 부동소수점 구간의 경계마다 정수와 실수를 오가는 변환 연산이 추가된다. 그 결과 지연이 FP32 대비
30\% 줄어드는 데 그쳤다. 반면 OpenVINO 경로는 헤드의 합성곱은 양자화하고 시그모이드, DFL, 원소별 산술 연산만
부동소수점으로 남겼다. 출력 붕괴를 일으키는 민감한 연산만 골라 제외하고 연산량의 대부분은 정수로 옮긴 것이다.
OpenVINO 런타임이 CPU의 정수 벡터 명령에 맞춰 연산을 최적화한다는 점도 기여했을 것이나, 그 몫을 분리하지는 않았다.

정확도 측면에서도 OpenVINO 경로가 나았다(Dice 0.8290 대 0.8184). 두 경로 모두 가중치는 대칭, 활성값은 비대칭으로 양자화하므로
기본 방식은 비슷하다. 그러나 캘리브레이션 데이터(TN3K 학습·검증 150장 대 fold별 검증 분할), 양자화 제외 범위, NNCF가 기본으로
수행하는 편향 보정 등 여러 요인이 함께 다르므로, 정확도 차이를 한 요인으로 돌리기는 어렵다. 이 비교에서 얻을 수 있는 실무적 교훈은, 같은 'INT8 사후 양자화'라는
이름 아래서도 어떤 연산을 어떻게 제외하느냐가 결과를 결정한다는 것이다.

\section{연구의 한계}\label{sec:limitations}

본 연구의 결과는 다음 한계 안에서 해석해야 한다.

\begin{enumerate}
  \item \textbf{계수 선택 편향.} 경계 손실 계수는 held-out 집합의 결과로 선택되었다(\secref{sec:sweep}). 사후 점검에서
    경계 정교화의 효과 자체는 검증 분할에서도 재현되었지만, 채택한 구성의 우위는 재현되지 않았다. 따라서 최종 모델의
    개선 폭은 낙관적일 수 있다. 검증 분할에서 가장 나았던 구성($g=0.25$, $\lambda=2$)을 전체 fold에서 다시 학습하여
    held-out에서 평가하면 이 편향 없이 효과를 추정할 수 있다.
  \item \textbf{작은 효과크기.} 경계 개선의 효과크기는 $d_z\approx0.1$--0.15로 작다. 통계적으로 유의한 차이가
    판독에 의미 있는 차이인지는 검증하지 않았다.
  \item \textbf{교차 데이터셋 결과의 낮은 검정력.} 시드 3개로는 작은 효과를 검출하기 어렵다. 교차 데이터셋에서
    개선도 악화도 확인되지 않았다는 결과는 '효과가 없다'가 아니라 '판정할 수 없다'로 읽어야 한다.
  \item \textbf{LODO의 낙관적 평가.} LODO에서는 타깃 데이터셋을 검증 집합으로도 사용하여 체크포인트를 골랐으므로,
    두 모델의 비교는 공정하지만 절대값은 실제 미지 도메인의 성능보다 높을 수 있다.
  \item \textbf{학습 무작위성의 부분적 반영.} in-domain 실험은 fold마다 시드 42로 한 번씩 학습하였다. fold 간 변동이
    학습 무작위성의 일부를 반영하지만, 같은 fold를 여러 시드로 반복하지는 않았다.
  \item \textbf{지연 측정 환경.} CPU 지연은 96코어 서버 CPU(AMD EPYC 9654)에서 측정하였다. 휴대형 장비나 일반 판독 PC의
    저전력 CPU에서는 절대값이 크게 다를 것이며, 정밀도 간 상대 순위도 달라질 수 있다.
  \item \textbf{범위의 제한.} 하나의 구조(YOLOv11n-seg), 단일 병변 클래스, 두 공개 데이터셋에 한정된 결과이다.
    임상 기관에서 새로 수집한 외부 데이터로는 검증하지 않았다.
\end{enumerate}
